\section{A weaker unbounded count cone with an earlier all-padding C entrance}
For an actual normalized arm $1^\alpha2^uV$, let $V$ have $r\ge1$
capacities at least three, $h=\sum(t-1)$, $W=h+r$, $w=u+r$, and
$\delta=h-u$. We prove every C shift at least three at every current
$N\ge W$ under
\[
               4w+30\ge13\delta,
       \quad\hbox{equivalently }17u+4r+30\ge13h.       \tag{1}
\]
This enlarges the registered cone with constant seventeen and moves
its entrance from $W+1$ to $W$. General earlier C, full A2 and the
original unamplified coefficient-unimodality conjecture remain open.

The one-factor family has registered complete C. The case
$\delta\le4$ uses the registered all-padding halfline from $W-1$;
$W\le32$ uses the registered complete nonbinary-weight32 theorem.
Hence assume $r\ge2$, $\delta\ge5$, $W\ge33$, and put $n=\delta-2$.
All parameters in this remaining case are unbounded.

\subsection*{The two CI head growth rates}
Put $q=\phi-1$, so $q^2+q=1$, and set
\[
 L=\prod[t],\quad E=(1+z)^uL,\quad
 f_j=[z^j]E/(1+qz),\quad \rho=q^h|L(-\phi)|.
\]
There are exactly $w$ interval factors in $E$. The registered
coefficientwise differential inequality gives
$jE_j\ge(w-j+1)E_{j-1}$.
Cone (1) gives $4w+4\ge13n$, so every ratio through $n$ is at
least $\gamma=9/4$. It also gives
\[
                 8(w-n)\ge15(n+1).                 \tag{2}
\]
For $n\ge8$ use $4w\ge13n-4$; for $n\le7$ use $w+n+2=W\ge33$.
Consequently $E_{n+1}\ge(15/8)E_n$. These are exhaustive integer
alternatives. In particular the entire head through $n+1$ is
nondecreasing. The registered divided-head induction therefore gives
$0<f_j\le E_j$, $f_j=E_j-qf_{j-1}$ at all needed indices.
The equivalent symmetric-middle guard $h\le3u+2$ also follows from
(1), $2r\le h$ and $\delta\ge5$; it is kernel checked.

\subsection*{A stronger resource from the existing reference amplification}
Let $R$ be the same-tail infinite-BINARY reference, $g$ its cumulative,
$x=D^{\rm ref}_W$, $y=D^{\rm ref}_{W+1}$, $P=\prod t$, and let $M$
be its numerator mass. Registered exact identities give $0<M<P$,
$y<\phi x$, $y\ge8x/5$, and
\[
 g_{W-4}=x-M+b_1,\quad g_{W-3}=y-M+1,
 \quad D^{\rm ref}_{W-1}=y-x+1,
\]
where $b_1=r-\mathbf1_{t_1=3}\ge1$.
The registered all-capacity $8/7$ amplification and positive-conjugate
Binet formula imply
$x\ge P(8/7)^W/(\phi\sqrt5)$.
For every $W\ge17$, $(8/7)^W\ge585/64$, whereas
$\phi\sqrt5\le117/32$. Hence $x\ge(5/2)P$ and
\[
                            M\le2x/5.              \tag{3}
\]
The inequality $P\ge W+3$ for two or more intervals of length at least
three follows by a full finite-product induction. Thus $x\ge2P\ge40$.
Using $q\le5/8$ gives
\[
 D^{\rm ref}_{W-1}\le(13/20)x,\quad
 g_{W-4}\ge3x/5,\quad g_{W-3}\ge6x/5.               \tag{4}
\]
The reference cumulative has the registered dyadic UPPER growth
$g_j\le2g_{j-1}$ and golden LOWER ratios $g_j>\phi g_{j-1}$.
No cumulative LC of the actual arm is assumed.

\subsection*{Exact finite padding increases the head golden gap}
Write $\widetilde D_N=F_N-F_{N-1}$ for signed ordinary differences and
$D_N=\max(\widetilde D_N,0)$. Infinite INITIAL padding here keeps $u$
fixed. The exact physical inverse, including its previous numerator
unit, gives
\[
 \phi D^\infty_{W-1}-D^\infty_W=C_0+f_n,
 \quad \phi D^\infty_W-D^\infty_{W+1}=\phi\rho+f_{n+1},
 \quad C_0=\phi-\phi^2\rho\ge q^5>0.                \tag{5}
\]
Finite padding subtracts its shifted $E$ coefficient from each signed
derivative. Therefore its golden correction at any read is
\[
                         E_j-\phi E_{j-1},          \tag{6}
\]
with negative coefficient indices zero. This sign is essential.
All head ratios read in (6) exceed $\phi$, including the last
$15/8$ ratio, so every correction is nonnegative.
From $f_n\ge E_n-qE_{n-1}$ and $q\le5/8$ one obtains the uniform
physical-$W$ bound
\[
 \phi\widetilde D_{W-1}-\widetilde D_W
                         \ge(13/8)E_{n-1}.         \tag{7}
\]
At current $W+1$, put $a=D_W$, $b=D_{W+1}$, $\Gamma=\phi a-b$ on
the positive current pair. With $E_p=E_{n-1}$, $E_c=E_n$, the exact
next-head bound and (6) give
\[
 \Gamma\ge(5/4)E_c\quad\hbox{for every padding}.     \tag{8}
\]
At $\alpha=1$ its extra term is $E_{n+1}-\phi E_n\ge E_n/4$;
at $\alpha=2$ it is $E_n-\phi E_{n-1}\ge5E_{n-1}/8$. Hence
\[
 \alpha=1:\ \Gamma\ge(3/2)E_c,\qquad
 \alpha=2:\ \Gamma\ge(5/4)E_c+(5/8)E_p.             \tag{9}
\]
No relative-LC hypothesis for finite/infinite padding is used.

\subsection*{Two complete prefix inductions bound the full loss}
The exact derivative loss relative to $R$ is
$\ell=E(H_2+z^{\alpha-1})$, with $H_2=(1-z-z^2)^{-1}$.
Let $S$ be its cumulative and $T$ the cumulative of $S$.
The response $H=EH_2$ satisfies
$H_j=E_j+H_{j-1}+H_{j-2}$.
Induction over every needed prefix, using $E_{j-1}\le4E_j/9$ and
$E_{j-2}\le16E_j/81$, proves $H_j\le81E_j/29$.
The shifted $E$ term is at most $E_j$, so $\ell_j\le110E_j/29$.
A second full prefix induction in $S_j=S_{j-1}+\ell_j$ proves
\[
                            S_j\le198E_j/29.        \tag{10}
\]
Both inductions are kernel checked for arbitrary prefix length. If a
current positive derivative pair is increasing, $\lambda=b/a>1$,
then $T_j-\lambda T_{j-1}\le S_j$.

\subsection*{Pay every read at current $W$ and the C reads at $W+1$}
At current $W$, every unprotected C read is $k=W-d$ with $d\ge3$,
$j=k-(w+2)=n-d\ge0$. Its previous actual derivative is no larger
than $D^{\rm ref}_{W-1}$. From (4), dyadic upper growth, and (7),
\[
 { (\phi D_{W-1}-D_W)g_{k-1}\over D_{W-1}}
 \ge {13\over8}{12\over13}
       {\gamma^{d-1}\over2^{d-3}}E_j
 \ge {243\over32}E_j>{198\over29}E_j.                \tag{11}
\]
Here $E_{n-1}\ge\gamma^{d-1}E_j$ by an all-interval growth induction;
the distance multiplier increases by $9/8$. The exact last difference
in (11) is $711/928>0$. Equation (10) pays the full double cumulative
loss, while the reference golden ratio adds another nonnegative term.
Protected reads are immediate. Decreasing/clipped current pairs are
also immediate by nonnegative cumulatives and actual slope support.
This proves every C shift at current $W$.

For current $W+1$, write $k=W+1-d$, $d\ge3$, so $j=n+1-d$.
Now $a\le x$ and $g_{W-3}\ge6x/5$. Replacing the two first factors
in (11) by $(5/4)(6/5)$, and using (8), gives the SAME $243/32$
payment. Thus every actual $G_k\ge\lambda G_{k-1}$ through $k\le W-2$
is established. The last two strong-prefix reads require the next step.

\subsection*{Two source-paid ordinary seeds complete the entrance}
Put $e_0=\phi F_{W-1}-F_W$ and $e_p=\phi F_{W-2}-F_{W-1}$.
The exact infinite-padding formulas are
\[
 e_0^\infty=U_n-\rho,\qquad
 e_p^\infty=U_{n-1}-\phi(1-\rho),\quad
 U_j=q\sum_{i\le j}E_i+q^2f_j.
\]
CI geometric growth gives $\sum_{i\le j}E_i\le9E_j/5$, hence
$U_j\le3E_j/2$. This prefix induction and kernel bound are formalized.
The ordinary finite-padding correction is $E_j-q\sum_{i<j}E_i$,
with indices $j=n+1-\alpha$ and $n-\alpha$ at the two reads.
For $\alpha=1$ the cumulative terms cancel exactly; for the other
paddings use the appropriate one- or two-step CI ratio. This gives
\[
\begin{array}{c|cc}
 &e_0&e_p\\\hline
\alpha=1&2E_c&2E_p\\
\alpha=2&(35/18)E_c&(35/18)E_p\\
\alpha\ge3&(35/18)E_c&(275/162)E_p
\end{array}                                                     \tag{12}
\]
as upper bounds.

Only $b>a>0$ requires an entrance. The registered actual ordinary
golden ceiling at $W$ gives $F_{W-1}\ge\phi a$. Actual positive-D
unimodality gives $D_{W-1}\le a$ in this strictly growing-entry case,
so $F_{W-2}\ge q a$. Use $\phi\ge8/5$, $q\ge8/13$ and
$E_c\ge9E_p/4$ in (8)--(12). In all three padding cases,
\[
                         \phi\Gamma>e_0,
                           \qquad q\Gamma>e_p.       \tag{13}
\]
For example the smallest previous-seed margin, in the third case,
is $(45/26-275/162)E_p=(35/1053)E_p>0$.
Thus both actual ordinary determinants are positive:
\[
 aF_W-bF_{W-1}=F_{W-1}\Gamma-ae_0>0,
\]
\[
 aF_{W-1}-bF_{W-2}=F_{W-2}\Gamma-ae_p>0.             \tag{14}
\]
The paid $G$ comparison at $W-2$ and the second inequality in (14)
give $G_{W-1}\ge\lambda G_{W-2}$; the first then gives
$G_W\ge\lambda G_{W-1}$. Every strong-prefix read through $W$,
and the ordinary seed at $W$, are now paid.

The registered universal pre-first-mode D-LC integrates these seeds
to every later C comparison. At a later strictly growing current use
its OWN ratio $\lambda_N\le\lambda$, with every earlier positive
tail ratio at least $\lambda_N$. The difference
$F_j-\lambda_NF_{j-1}$ is then nondecreasing from its positive seed;
the same holds after one further summation for $G$. Earlier prefix
reads were already paid with the larger entrance ratio. On and after
the first D mode C is immediate. If the entrance pair is flat,
pre-mode D-LC excludes a later increase; ordinary D-UI alone is not
used to exclude a pre-mode flat. This completes the entire halfline.

\subsection*{Uniform consequences and exact scope}
If $\delta\le9$ and (1) fails, then $W\le30$. The old weight32
theorem therefore gives, with (1), EVERY actual C comparison at every
$N\ge W$ for ALL $u\ge h-9$. No new finite entrance is introduced.
The registered all-shift boundary at $W-1$ when $u\ge h-5$ now gives
the continuous C halfline $N\ge W-1$ at ALL $u\ge h-5$. This closes
the previously retained $u=h-5$, current-$W$ gap.

Two standalone Lean files check twenty-eight scoped reports, including
complete arbitrary-prefix renewal, cumulative-loss, CI-prefix and
geometric-growth inductions, a finite-product induction, all three
padding budgets, full read-payment algebra and all-integer cone guards.
The exact actual/reference/root/CI source identifications and complete
sequence/mode assemblies are the written proof above; there is no
end-to-end Lean theorem for actual arms. Every new finite profile is
QA only; the old reference exception certificates and weight32 theorem
are reused with their previous scope and are not recounted.
General earlier C, the complementary count region, full A2 and original
unamplified coefficient unimodality remain OPEN.
